\documentclass[10pt]{article}

% Portable arXiv-compatible Times-family typography.  The Microsoft
% Times New Roman font is not bundled with TeX Live; mathptmx supplies
% a metrically familiar Times face and matching mathematics.
\usepackage[T1]{fontenc}
\usepackage[utf8]{inputenc}
\usepackage{mathptmx}
\usepackage{microtype}
\usepackage[
  letterpaper,
  left=1in,
  right=1in,
  top=0.72in,
  bottom=0.68in
]{geometry}
\usepackage{mathtools,amssymb,amsthm}
\usepackage{booktabs}
\usepackage{enumitem}
\usepackage{xcolor}
\usepackage{fancyhdr}
\usepackage{titlesec}
\usepackage[
  colorlinks=true,
  linkcolor=blue!55!black,
  citecolor=blue!55!black,
  urlcolor=blue!55!black
]{hyperref}
\hypersetup{
  pdftitle={A Hamiltonian-Path Criterion from Bipartite Number and Average Eccentricity},
  pdfauthor={Luke Kabbash},
  pdfsubject={Graph theory; Written on the Wall II, Conjecture 198a},
  pdfkeywords={Hamiltonian path, induced bipartite subgraph, eccentricity, graph radius}
}

\definecolor{paperink}{HTML}{171717}
\definecolor{papermuted}{HTML}{565B60}
\definecolor{paperaccent}{HTML}{294C69}
\definecolor{paperhairline}{HTML}{AEB5BB}
\color{paperink}
\hypersetup{
  linkcolor=paperaccent,
  citecolor=paperaccent,
  urlcolor=paperaccent
}

\setlength{\parindent}{0pt}
\setlength{\parskip}{6.4pt}
\linespread{1.04}
\setlist{topsep=2pt,partopsep=0pt}

\titleformat{\section}[block]
  {\normalfont\bfseries\fontsize{13.5}{16.4}\selectfont\color{paperink}}
  {\thesection.}{0.45em}{}
\titlespacing*{\section}{0pt}{12pt}{6pt}
\titleformat{name=\section,numberless}[block]
  {\normalfont\bfseries\fontsize{13.5}{16.4}\selectfont\color{paperink}}
  {}{0pt}{}
\titleformat{\subsection}[block]
  {\normalfont\itshape\fontsize{11.2}{14}\selectfont\color{paperink}}
  {\thesubsection}{0.45em}{}
\titlespacing*{\subsection}{0pt}{8pt}{4.5pt}

\pagestyle{fancy}
\fancyhf{}
\fancyhead[L]{%
  \sffamily\fontsize{7}{8}\selectfont\color{papermuted}%
  HAMILTONIAN-PATH CRITERION}
\fancyhead[R]{%
  \sffamily\fontsize{7}{8}\selectfont\color{papermuted}\thepage}
\renewcommand{\headrulewidth}{0.35pt}
\renewcommand{\headrule}{%
  \hbox to\headwidth{\color{paperhairline}\leaders\hrule height
  \headrulewidth\hfill}}
\setlength{\headheight}{17.6pt}
\fancypagestyle{plain}{%
  \fancyhf{}
  \fancyfoot[C]{%
    \sffamily\fontsize{7}{8}\selectfont\color{papermuted}\thepage}
  \renewcommand{\headrulewidth}{0pt}
}

\renewenvironment{abstract}
  {\begin{list}{}{%
     \setlength{\leftmargin}{0.22in}%
     \setlength{\rightmargin}{0.22in}%
   }\item\relax}
  {\end{list}}

\IfFileExists{artifact-status.tex}
  {% The exact repository theorem and its radius bound are kernel-checked.
\newif\ifleanverified
\leanverifiedtrue
\newcommand{\leanstatus}{Lean 4.27 kernel checks pass for the exact
repository theorem, including
\(b(G)\geq 2\,\operatorname{rad}(G)\), with no proof placeholders.}
}
  {\newif\ifleanverified\leanverifiedfalse
   \newcommand{\leanstatus}{Formal verification status was not supplied.}}

\newtheoremstyle{paperplain}
  {0pt}
  {0pt}
  {\itshape}
  {0pt}
  {\bfseries}
  {.}
  {\newline}
  {}
\theoremstyle{paperplain}
\newtheorem{theorem}{Theorem}
\newtheorem{lemma}[theorem]{Lemma}
\newtheorem{remark}[theorem]{Remark}
\newcommand{\ecc}{\operatorname{ecc}}
\newcommand{\rad}{\operatorname{rad}}
\newcommand{\diam}{\operatorname{diam}}
\newcommand{\avg}{\operatorname{avg}}

\makeatletter
\renewcommand{\maketitle}{%
  \thispagestyle{plain}
  \begin{center}
    \vspace*{6pt}
    {\bfseries\fontsize{23.5}{26.5}\selectfont
      A Hamiltonian-Path Criterion from\\
      Bipartite Number and Average Eccentricity\par}
    \vspace{7pt}
    {\itshape\fontsize{11.2}{14.2}\selectfont\color{papermuted}
      A proof of Written on the Wall II, Conjecture 198a\par}
    \vspace{7pt}
    {\fontsize{10.2}{13}\selectfont
      Luke Kabbash \textbar{} GPT-5.6 Sol\par}
    \vspace{1pt}
    {\fontsize{9.2}{12}\selectfont\color{papermuted}23 July 2026\par}
  \end{center}
  \vspace{10pt}
  {\color{paperink}\hrule height 0.75pt}
  \vspace{10pt}
}
\makeatother

\begin{document}
\maketitle

\begin{abstract}
\small
\noindent\textbf{Abstract.}\enspace
For a finite nontrivial connected simple graph \(G\), let \(b(G)\) be
the maximum order of an induced bipartite subgraph and let
\(\overline{\ecc}(G)\) be the average vertex eccentricity.  We prove
that
\[
        b(G)\leq 2+\overline{\ecc}(G)
        \quad\Longrightarrow\quad
        G\ \text{has a Hamiltonian path}.
\]
This is the statement recorded as Written on the Wall II,
Conjecture~198a.  The proof is driven by the tight chain
\(D+1\leq b(G)\leq 2+\overline{\ecc}(G)\leq D+2\), where
\(D=\diam(G)\).  In the first equality case, every vertex outside a
diametral path belongs to a clique layer that can be inserted into
that path.  In the second, the graph is self-centered of diameter two,
and the Chv\'atal--Erd\H{o}s traceability theorem applies.
\end{abstract}

\section{Statement and notation}

All graphs in this note are finite and simple.  For a connected graph
\(G\), write
\[
  D=\diam(G),\qquad
  r=\rad(G),\qquad
  \overline{\ecc}(G)
    =\frac{1}{|V(G)|}\sum_{v\in V(G)}\ecc_G(v).
\]
The \emph{bipartite number} \(b(G)\) is the largest cardinality of a
set \(S\subseteq V(G)\) for which the induced graph \(G[S]\) is
bipartite.  A graph is \emph{traceable} if it has a Hamiltonian path.

\vspace{1pt}
\noindent{\color{paperink}\rule{\linewidth}{0.75pt}}
\vspace{-2pt}

\begin{theorem}[Written on the Wall II, Conjecture 198a]
\label{thm:main}
Let \(G\) be a nontrivial connected graph.  If
\[
                   b(G)\leq 2+\overline{\ecc}(G),
\]
then \(G\) is traceable.
\end{theorem}

\vspace{-3pt}
\noindent{\color{paperink}\rule{\linewidth}{0.75pt}}
\vspace{2pt}

We use two classical facts.  First,
\begin{equation}
                         b(G)\geq 2r ,
\label{eq:radius-bipartite}
\end{equation}
as recorded by DeLaVi\~na, Pepper, and Waller
\cite[Theorem~2]{DeLaVinaPepperWaller2007}.  Second, the
Chv\'atal--Erd\H{o}s path theorem says that an \(s\)-connected graph
with no independent set of \(s+2\) vertices is traceable
\cite[Theorem~2]{ChvatalErdos1972}.

\clearpage
\section{Proof}

Choose a diametral geodesic
\[
                         P=v_0v_1\cdots v_D .
\]
A geodesic is induced, so \(P\) is an induced bipartite graph of order
\(D+1\).  Since every eccentricity is at most \(D\),
\begin{equation}
       D+1\leq b(G)
       \leq 2+\overline{\ecc}(G)
       \leq D+2.
\label{eq:chain}
\end{equation}
Because \(b(G)\) is integral, only two cases remain.

\subsection*{Case 1: \(b(G)=D+1\)}

Fix \(x\notin V(P)\).  The induced graph on \(V(P)\cup\{x\}\) has
order \(D+2>b(G)\), hence is not bipartite.  Thus \(x\) has neighbors
on \(P\) of both parities.  On the other hand, any two \(P\)-neighbors
\(v_i,v_k\) of \(x\) satisfy \(|i-k|\leq2\): otherwise
\(v_i x v_k\) would shorten the corresponding section of the
geodesic \(P\).

Let \(j\) be the least index for which \(xv_j\in E(G)\).  The preceding
paragraph gives
\begin{equation}
       N_G(x)\cap V(P)
       \in\bigl\{\{v_j,v_{j+1}\},
                   \{v_j,v_{j+1},v_{j+2}\}\bigr\},
\label{eq:attachment}
\end{equation}
where the second alternative is omitted when \(j+2>D\).  For
\(0\leq j<D\), let \(X_j\) consist of the vertices outside \(P\) whose
least neighbor index is \(j\).

Each \(X_j\) is a clique.  Indeed, if distinct \(x,y\in X_j\) were
nonadjacent, then
\[
       \bigl(V(P)\setminus\{v_{j+1}\}\bigr)\cup\{x,y\}
\]
would induce a bipartite graph of order \(D+2\).  To see the
bipartition directly, retain the alternating coloring of \(P\):
the only remaining possible \(P\)-neighbors of \(x\) or \(y\) are
\(v_j\) and \(v_{j+2}\), which have the same color, and assign both
\(x\) and \(y\) the opposite color.  This contradicts \(b(G)=D+1\).

Now order the vertices as
\begin{equation}
            v_0,\ X_0,\ v_1,\ X_1,\ \ldots,\
            v_{D-1},\ X_{D-1},\ v_D.
\label{eq:ham-order}
\end{equation}
using any order inside each \(X_j\) and simply omitting empty layers.
Consecutive vertices are adjacent: every \(X_j\) is a clique, and
every member of \(X_j\) is adjacent to both \(v_j\) and \(v_{j+1}\).
All vertices occur exactly once, so \eqref{eq:ham-order} is a Hamiltonian
path.

\subsection*{Case 2: \(b(G)=D+2\)}

Equality throughout the right half of \eqref{eq:chain} gives
\(\overline{\ecc}(G)=D\).  Every vertex therefore has eccentricity
\(D\), so \(G\) is self-centered and \(r=D\).  Combining this with
\eqref{eq:radius-bipartite} yields
\[
                         2D=2r\leq b(G)=D+2,
\]
and hence \(D\leq2\).  The case \(D=1\) is impossible: a nontrivial
diameter-one graph is complete and has bipartite number \(2\), not
\(D+2=3\).  Consequently,
\begin{equation}
                         D=r=2,\qquad b(G)=4.
\label{eq:diam-two}
\end{equation}

The graph \(G\) has no cut vertex.  Otherwise, let \(c\) be a cut
vertex.  For any \(w\ne c\), choose \(y\) in a component of \(G-c\)
different from the component containing \(w\).  Every \(w\)-\(y\)
path passes through \(c\), and a path of length at most two must
therefore be \(wcy\).  Hence every vertex other than \(c\) is adjacent
to \(c\), giving \(\ecc_G(c)=1\), contrary to self-centeredness with
radius two.  Since \(D=2\) also implies \(|V(G)|\geq3\), the graph is
2-connected.

Finally, if \(I\) is a maximum independent set, then \(I\neq V(G)\)
because \(G\) is connected and nontrivial.
Adding any vertex outside \(I\) produces an induced bipartite graph,
so \(b(G)\geq |I|+1\).  By \eqref{eq:diam-two},
\(\alpha(G)\leq3\).  Hence this 2-connected graph has no independent
set of four vertices.  The Chv\'atal--Erd\H{o}s path theorem, with
\(s=2\), supplies a Hamiltonian path.  This completes the proof of
Theorem~\ref{thm:main}. \qed

\section{Formalization and reproducibility}

The theorem is represented in the Formal Conjectures repository over
a finite type \(\alpha\) by a connected simple graph \(G\), with
\(b(G)\) coerced to \(\mathbb{R}\) and average eccentricity defined as
a real average.  The proof artifact accompanying this manuscript is
organized around the following independently checkable interfaces:
\begin{enumerate}[leftmargin=2.1em,itemsep=2pt]
  \item a diametral geodesic and the lower bound \(D+1\leq b(G)\);
  \item the integral dichotomy \(b(G)\in\{D+1,D+2\}\);
  \item the attachment and clique-layer lemmas
        \eqref{eq:attachment};
  \item the self-centered diameter-two reduction; and
  \item a direct longest-path traceability endpoint matching the repository
        statement:
        \[
          \exists\,a\,b,\quad
          \exists p:G.\mathrm{Walk}\ a\ b,\quad
          p.\mathrm{IsHamiltonian}.
        \]
\end{enumerate}

All of these components, including both equality branches and the classical
inequality \eqref{eq:radius-bipartite}, are checked in Lean without proof
placeholders. Their top-level composition has the exact hypotheses and
conclusion of the repository conjecture.

\paragraph{Artifact status.}
\ifleanverified
\leanstatus
\else
\leanstatus\  This manuscript asserts the mathematical proof above,
but it does not by itself assert that a Lean kernel has accepted the
entire development.  The distributed verification log and checksum
manifest are authoritative for that separate claim.
\fi

\paragraph{Research provenance and generative-AI use disclosure.}
I used a near-autonomous, iteratively prompted research workflow to identify
and pursue this problem.  I began by examining other open problems reported
as solved during the preceding several days, comparing their statements and
proof strategies to assess which nearby conjectures might be approachable.
The resulting model outputs placed this conjecture among comparatively
steeper targets, and I selected it for sustained work.

I then used GPT-5.6 Sol in local Codex at extra-high and ultra reasoning
settings, cloud-hosted ChatGPT work, and individual ChatGPT 5.6 Pro sessions
for iterative research, proof development, Lean formalization, and
verification, continuing until the exact theorem passed the Lean kernel.
GPT-5.6 Sol at extra-high reasoning was used for repository preparation;
Claude Fable 5 and Claude Sonnet 5 were used for limited supplementary
searches.  Model names in the displayed byline and this disclosure identify
software used in the workflow, not independent agency, authorship, or
responsibility.  I am the sole human author and take responsibility for the
claims and submitted material.

\section*{Prior-solution search and research status}

The statement appears as Conjecture~198a in Written on the Wall II
\cite{DeLaVinaWOWII} and in the Formal Conjectures snapshot
\cite{FormalConjectures198a}.  A targeted search completed on 23 July
2026 covered the current upstream statement, its open and closed issue
and pull-request history, public GitHub code indexed under the exact
theorem identifier, and exact-phrase web queries.  It found copies of
the open conjecture and benchmark tasks, but no evidence in the searched
sources of a proof or verified formalization predating this work.  The
search was not an exhaustive literature review, does not establish
worldwide novelty or priority, and supports no claim of a first proof.

\bibliographystyle{alpha}
\bibliography{references}

\end{document}
